\section{Limitations, Objections, and Replies}
\label{sec:objections}

\noindent
A foundational treatment is strengthened by gathering, and meeting directly, the objections its apparatus invites, rather than leaving them dispersed among the cautions of the several sections. This section states the principal objections to the paper's approach and the replies the paper offers, in the awareness that the replies dissolve some objections and only mitigate others, and that the latter mark the genuine limitations the paper bears.

\subsection{That the skew criterion misjudges a legitimate division of labour}
\label{sec:obj-division}

The first objection holds that the skew of the return distribution will misjudge as exploitation a division of labour that is in fact legitimate. Parties to a shared life may, the objection runs, justly specialise, one bearing more of the sustaining labour and the other more of some other contribution, and a measure that reads asymmetric return as injustice will condemn a specialisation the parties have freely and fairly chosen.

The reply turns on the coupling-based baseline (\xref{sec:ev-return}). The criterion does not read asymmetric contribution as injustice, nor asymmetric return as such; it reads as injustice the systematic shortfall of one party's return below the baseline that the coupling and the generation set. A legitimate specialisation, in which one party bears more sustaining labour and is returned to in proportion, through the relation's circulation, to the value that labour generates, produces no skew, since the return tracks the generation; the normalised return of both parties is near unity, and the distribution is symmetric. The skew arises only when the value a party generates does not return to that party, which is not the mark of specialisation but of extraction. The objection conflates an asymmetric division of tasks, which the criterion does not condemn, with an asymmetric return on a symmetric generation, which it does. The reply, however, concedes a residual difficulty: the baseline itself depends on the assignment of value to contributions (\xref{sec:ev-operational}), and a culture or a couple that systematically undervalues the sustaining labour in that assignment will set a baseline that conceals the extraction, reading a genuine shortfall as proportionate. The criterion is only as just as the valuation that feeds it, and the valuation is the contestable step the concession marks.

\subsection{That modelling a relation as a value network reifies intimacy}
\label{sec:obj-reification}

The second objection holds that to model a shared life as a network of value generation and return is to reify intimacy, to reduce a relation of love to an economic exchange, and that the reduction falsifies the very thing it claims to illuminate.

The reply distinguishes the model from a metaphysics. The paper does not claim that a relation is a value network, nor that love is exchange; it claims that the network is a structured proxy (\xref{sec:bg-concession}), an operational instrument that makes one feature of a relation, the distribution of its returns, assessable, and that this feature, though it is not the relation, is a feature a relation has and one whose maldistribution is an injustice the parties have reason to detect. The value at issue is, throughout, the transversal and largely unsymbolised value of the prior papers (\xref{sec:appraisal}), not a monetary or exchangeable quantity, and the paper has insisted that the structured proxy does not reach it. The reduction the objection fears is exactly the reduction the operational concession refuses: to mistake the proxy for the relation, the measured return for the love, is the cardinal error the paper has named (\xref{sec:ref-possession}), and the paper commits it no more than a cartographer, mapping a country's rivers, claims that the country is its rivers. The model illuminates a feature against the background of a relation it does not claim to exhaust, and the warmth, the meaning, and the unsymbolised texture of the shared life, which the model does not capture, are precisely what the paper's Envoi and its insistence on the matter of the field return to.

\subsection{That the operational concession immunises the paper against criticism}
\label{sec:obj-concession}

The third objection holds that the operational concession is too convenient: by declaring in advance that its measures are proxies that do not reach the real value, the paper immunises itself against every criticism, since any objection to a measure can be met by the reply that the measure was never claimed to be more than a proxy. The concession, on this objection, is not a discipline but an evasion.

The reply grants the danger and specifies the discipline that averts it. A concession that licensed any measure whatever, on the ground that all are merely proxies, would indeed be an evasion; the paper's concession does not. It requires of a proxy that it be operationally warranted, that it make some genuine feature assessable, that its results be read as evidence and not as measurement, and that it never be possessed as the truth of the relation (\xref{sec:ref-possession}); and it holds the proxy answerable to the basal appraisal it proxies, so that a measure systematically at odds with the parties' considered appraisal is thereby impugned. The concession is, in this respect, a standard the measures must meet and not a blanket exemption: it constrains how a measure may be used, and the constraint is violable, as the cardinal error of possessing the measure as truth shows. The objection would hold against a concession that exempted measures from criticism; it does not hold against one that, as here, specifies the limited use under which alone a measure is admitted and marks the misuse that breaks it.

\subsection{That to speak of justice in intimacy is a category mistake}
\label{sec:obj-category}

The fourth objection holds that justice is a category proper to the public and political sphere, to the distribution of goods among strangers under institutions, and that to import it into intimacy is a category mistake, since intimate relations are governed by love and care, which transcend and even oppose the calculative posture of justice.

The reply is that the opposition between love and justice is itself the ideology under which intimate injustice is concealed, and that the ethics of care, properly understood, requires rather than opposes the justice of care's distribution (\xref{sec:phil-justice}). The labour of sustaining a shared life is labour; its systematic extraction from one party is an injustice whether or not the relation is loving, and indeed a loving relation is precisely the site at which such extraction is most effectively concealed, since the language of love supplies the ideology under which the sustaining labour is figured as freely given and its under-return as natural (\xref{sec:ev-skew}). To hold that justice has no place in intimacy is not to protect intimacy from a calculative intrusion but to remove from the exploited party the one vocabulary in which their exploitation could be named. The paper's claim is not that intimacy should be calculative, but that the refusal to assess the justice of an intimate relation's distribution is the condition under which its injustice is sustained, and that a love worthy of the name does not require, and is not damaged by, the justice of its returns. The reply, however, concedes a limit: the assessment of justice is an instrument to be drawn upon when injustice is suspected, not a posture to be maintained continuously, and a relation conducted as a perpetual audit of its returns would have made the measure the truth of the relation, committing the cardinal error in a new form. Justice is the internal condition of the good cycle (\xref{sec:reg-justice}), not the continuous activity of the parties within it.

\subsection{The residual limitations}
\label{sec:obj-residual}

The replies dissolve the objections that rest on misunderstanding, the conflation of asymmetric return with asymmetric task, the mistaking of the model for a metaphysics, the reading of the concession as exemption, and the opposition of love to justice. They leave standing a set of genuine limitations, which the paper records as its own. The valuation of contributions, on which the whole quantitative apparatus rests, is an operational stipulation that no measurement discharges, and a biased valuation yields a biased assessment (\xref{sec:obj-division}). The network model is unvalidated and its estimation an unsolved empirical problem (\xref{sec:app-justice}). The dynamical model is a structural analogy under the standing qualification that the relation's law is reconfigured by what it generates. The proxies for the unfolding of each party reach least adequately the value that lies nearest the unsymbolised. And the whole apparatus is offered for the central case of the dyad and its coupled network, with the cross-cultural variation (\xref{sec:crosscultural}) locating but not exhausting the forms the apparatus must take across the range of human shared lives. These are the limitations a foundational treatment bequeaths to the work it founds, and their statement is not a weakening of the foundation but a part of it, since a foundation that concealed its limits would be the false ground the series has throughout refused.
